\chapter{Nombres relatifs et fractions}
\label{corriges:02}
\section*{Exercice 2.1 -- Calcul exact}
La multiplication est prioritaire :
\[A=\frac76-\frac{10}{12}=\frac76-\frac56=\frac13.\]
\[B=-\frac35\times\frac{10}{9}=-\frac23.\]
Enfin $C=-9+9=0$ : les parenthèses changent la base de la puissance.

\section*{Exercice 2.2 -- Comparer et approcher}
On a $-7/12=-35/60$ et $-3/5=-36/60$. Donc $-7/12>-3/5$.
Comme $17/7=2{,}428571\ldots$, $2{,}42<17/7<2{,}43$ et l'arrondi au centième est $2{,}43$. Le chiffre des millièmes est $8$.

\section*{Exercice 2.3 -- Une réserve}
Après le premier usage, il reste $60(1-1/3)=40\unite L$. Après le second, il reste $40(1-3/8)=25\unite L$. La fraction initiale conservée est $25/60=5/12$. On ne soustrait pas directement $1/3+3/8$, car les deux fractions ne portent pas sur le même volume de référence.

\section*{Exercice 2.4 -- Prouver une règle}
Multiplier le membre de gauche par $b$ donne $a-c$ ; ce membre est donc l'unique quotient de $a-c$ par $b$. Pour la seconde affirmation, choisir $a=1,b=2,c=1$ donne $2/3\ne1/2$. Ajouter la même quantité au numérateur et au dénominateur ne conserve pas généralement une fraction.
